\documentclass[11pt,a4paper]{article}
\usepackage[utf8]{inputenc}
\usepackage[T1]{fontenc}
\usepackage[ngerman]{babel}
\usepackage{lmodern}
\usepackage{microtype}
\usepackage{csquotes}
\usepackage[margin=2.5cm]{geometry}
\usepackage{booktabs}
\usepackage{tabularx}
\usepackage{amsmath}
\usepackage{xcolor}
\usepackage{pgfplots}
\pgfplotsset{compat=1.17, /pgf/number format/1000 sep={}}
\usepackage[hidelinks]{hyperref}
\usepackage{caption}
\captionsetup{font=small,labelfont=bf}

\definecolor{blau}{HTML}{2A78D6}
\definecolor{rot}{HTML}{E34948}
\definecolor{grau}{HTML}{888780}

\title{Super-TikTakToe:\\ Ein selbstlernendes Netz für Ultimate Tic-Tac-Toe\\[0.3em]
\large Rückblick, Ergebnisse und Diskussion}
\author{Kaspar Hui}
\date{21.~November 2026}

\begin{document}
\maketitle

\begin{abstract}
\noindent
Dieses Papier fasst ein Freizeitprojekt zusammen: ein Programm, das Ultimate Tic-Tac-Toe nach dem Vorbild von
AlphaZero allein durch Spiele gegen sich selbst lernt, ohne menschliches Wissen und ohne Minimax im Spiel.
Nach einer Reihe von Versuchen zu Netzgröße, Netzaufbau, Trainingseinstellungen und Suche zeigte sich ein
klares Muster: Fast alle Änderungen endeten auf demselben Niveau, und nur \emph{mehr Partien} brachten
messbar mehr Stärke, sofern das Netz groß genug war. Das finale Netz (1{,}84 Millionen Gewichte) wurde
deshalb mit gut einer Million zusätzlicher Partien weitertrainiert. In Endspielen auf 300 verschiedenen
Eröffnungen erreicht es gegen das bisher beste Netz (\emph{large 240}) 61{,}3\,\% bei 500 und 58{,}2\,\% bei
2000 Simulationen pro Zug und ist nur in 22 bzw.\ 19 von 300 Eröffnungen schlechter. Gegen eine
Minimax-Suche der Tiefe 15 gewinnt es mit 2000 Simulationen 88{,}2\,\% der Punkte und ist nur in einer
von 102 Eröffnungen schlechter.
\end{abstract}

\section{Einleitung}

\paragraph{Das Spiel.} Ultimate Tic-Tac-Toe besteht aus neun kleinen Tic-Tac-Toe-Feldern, die zusammen ein
großes Feld bilden. Wer ein kleines Feld gewinnt, besetzt das entsprechende Feld des großen Bretts; wer
drei davon in einer Reihe hat, gewinnt. Die Besonderheit: Das Feld, in das man setzt, bestimmt, in
welchem kleinen Feld der Gegner als Nächstes spielen muss. Dadurch hängt jeder Zug von der ganzen
Stellung ab, und kurze Rechnungen reichen selten. Für die hier verwendeten Regeln ist das Spiel nicht
gelöst.

\paragraph{Vorgeschichte.} Das Spiel selbst und eine erste Computergegnerin entstanden im November und
Dezember 2023: eine Minimax-Suche mit handgeschriebener Bewertung, die je nach Tiefe sehr stark spielt,
aber pro Zug lange rechnet. Sie blieb im Projekt als \emph{Messgegner}: Gegen Minimax der Tiefen 12 bis 15
lässt sich die Stärke eines Netzes von außen prüfen.

\paragraph{Ziel.} Im November 2026 kam der Teil hinzu, um den es hier geht: ein neuronales Netz, das nach
dem Vorbild von AlphaZero~\cite{alphazero} nur aus eigenen Partien lernt. Zwei Vorgaben prägten die Arbeit.
Erstens ein \emph{reines Netz}: keine Minimax-Hilfe im Spiel oder im Training. Zweitens ein Ziel, das über
die Punktzahl hinausgeht: Das finale Netz soll \emph{wenige Fehler} machen. Ein Ergebnis von 50\,\%, bei
dem beide Seiten auf denselben Eröffnungen mal gewinnen und mal verlieren, heißt, dass beide Fehler
machen. Das Netz soll auf möglichst wenigen Eröffnungen schlechter spielen als sein Gegner.

\section{Methode}

\subsection{Lernen durch Selbstspiel}

Das Verfahren folgt AlphaZero. Ein Netz schaut auf eine Stellung (sechs Bildebenen von $9\times9$
Feldern: eigene Steine, gegnerische Steine, erlaubte Züge, entschiedene Felder und Ähnliches) und gibt
zwei Dinge zurück: eine Verteilung über die 81 möglichen Züge (\emph{Policy}) und eine Einschätzung, wie
die Partie ausgeht (\emph{Value}, von $-1$ bis $+1$). Eine Monte-Carlo-Baumsuche (PUCT, $c_\text{puct}=1{,}5$)
fragt das Netz in jeder Simulation und wählt am Ende den meistbesuchten Zug.

Jede Trainingsrunde besteht aus zwei Teilen: Zuerst spielt das aktuelle Netz viele Partien gegen sich
selbst (512 bis 1024 gleichzeitig). Die Stellungen kommen mit den Besuchszahlen der Suche und dem
Ausgang der Partie in einen Speicher. Dann lernt das Netz aus Stichproben dieses Speichers, die
Besuchszahlen und den Ausgang vorherzusagen. Jede Stichprobe wird zufällig mit einer der acht
Symmetrien des Bretts (Drehen, Spiegeln) gedreht. Die Lernrate sinkt in Stufen.

\subsection{Netze}

Alle Netze sind Residual-Netze aus Faltungsschichten in drei Größen:

\begin{center}
\begin{tabular}{lrrr}
\toprule
Netz & Kanäle & Blöcke & Gewichte \\
\midrule
small  & 32 & 4 & 0{,}19 Mio. \\
medium & 64 & 6 & 0{,}67 Mio. \\
large  & 96 & 8 & 1{,}84 Mio. \\
\bottomrule
\end{tabular}
\end{center}

\noindent Trainiert wurde auf Google Colab (Grafikkarte L4), Endspiele und Auswertungen liefen auf einem
MacBook Air.

\subsection{Messen}

Gemessen wurde auf drei Arten:

\begin{enumerate}
\item \textbf{Checks während des Trainings:} alle 4 bzw.\ 8 Runden 256 oder 512 Partien gegen ein festes
Vergleichsnetz mit 100 Simulationen pro Zug. Schnell, aber sie überschätzen den Unterschied bei längerer
Bedenkzeit um 2 bis 4 Punkte.
\item \textbf{Direkte Partien} zwischen fertigen Netzen mit 500 und 2000 Simulationen, jede Eröffnung
einmal mit jeder Farbe. Ohne Zufall spielt dasselbe Paar aus derselben Stellung immer dieselbe Partie.
\item \textbf{Partien gegen Minimax} der Tiefen 13 und 15.
\end{enumerate}

\paragraph{Die Eröffnungsregel.} Für das Ziel \enquote{wenige Fehler} werden pro Eröffnung die Punkte beider
Partien addiert (Sieg 1, Unentschieden $\tfrac12$). Unter 1 heißt: auf dieser Eröffnung schlechter
gespielt als der Gegner; genau 1: gleich; über 1: besser. Diese Zählung trennt \enquote{viele Remis} von
\enquote{gleich viele Siege und Niederlagen}, die beide 50\,\% ergeben können.

\paragraph{Eröffnungen ohne Wiederholung.} Das Brett hat acht Symmetrien. Zwei Eröffnungen, die durch
Drehen oder Spiegeln ineinander übergehen, sind dieselbe Stellung. Die früher verwendeten 128 zufälligen
Eröffnungen enthielten nur 71 verschiedene Stellungen. Für die Endspiele wurden daher alle 102
verschiedenen Eröffnungen mit zwei Zügen und 198 der 822 verschiedenen mit drei Zügen verwendet, zusammen
300, jede von beiden Seiten gespielt.

\paragraph{Zufall.} Zwei Trainingsläufe mit identischen Einstellungen (\emph{large} und \emph{large\_repeat})
lagen nach 100 Runden etwa 3 Punkte auseinander. Ein Unterschied zwischen zwei Varianten zählt deshalb
erst ab etwa 4 Punkten über viele Checks; die Schwelle wurde jeweils vor einem Lauf festgelegt.

\section{Die Versuche vor dem finalen Lauf}

Tabelle~\ref{tab:versuche} fasst die Versuche zusammen. Gemeinsam ist fast allen: Sie endeten auf etwa
demselben Niveau.

\begin{table}[ht]
\centering
\small
\begin{tabularx}{\textwidth}{>{\raggedright\arraybackslash}p{4.6cm}X}
\toprule
Versuch & Ergebnis \\
\midrule
Größeres Netz (large statt medium), 240 Runden, 123\,000 Partien
 & Gegen medium 240 direkt nur 52{,}3\,\% (500 Sim.) bzw.\ 50{,}0\,\% (2000 Sim.). Die Größe allein half nicht. \\
Netzaufbau (acht Varianten D0--D8, gelernt aus den Stellungen von large)
 & \enquote{Linien-Schichten} (Summen entlang der 72 kleinen und 8 großen Linien) lernen etwa doppelt so schnell pro Lektion. \\
Linien-Netz selbst trainiert (\emph{large\_lines})
 & Führte in den Checks gegen medium (56{,}0\,\% gegen 51{,}1\,\%), verlor aber direkt gegen large (46{,}9\,\% / 47{,}7\,\%) und war 1{,}4- bis 1{,}5-mal langsamer. \\
Lektionen über den langen Spielverlauf (Besitzer der Felder, Ausgang als drei Wahrscheinlichkeiten, Zugwerte)
 & 0 bis 3 Punkte, innerhalb des Zufalls. \\
Kleines Netz mit vielen Partien (\emph{small\_many}, 215\,000 Partien)
 & Pro Stunde etwa gleichauf mit large, pro Partie weit zurück, flacht tiefer ab. Nach Stopp-Regel abgebrochen (43{,}7\,\% statt nötiger 48\,\%). \\
Doppelte Simulationen (800 statt 400)
 & Kein messbarer Gewinn (52{,}8\,\% gegen 52{,}7\,\%). \\
Doppelt so viele Partien pro Runde (\emph{more\_games})
 & $+2{,}4$ Punkte, selbst bei der kleinsten Lernrate. Der einzige Hebel, der deutlich half. \\
Suchparameter $c_\text{puct}$
 & Kleinere Werte helfen nur bei kurzer Bedenkzeit; 1{,}5 blieb. \\
Vergleich mit dem Netz von uttt.ai
 & medium spielte mit gleicher Suche ausgeglichen (47{,}0\,\% / 50{,}6\,\%). \\
\bottomrule
\end{tabularx}
\caption{Versuche vor dem finalen Lauf (Sim. = Simulationen pro Zug).}
\label{tab:versuche}
\end{table}

\paragraph{Lesart.} medium und large endeten bei gleich vielen Partien gleichauf; die Linien-Schichten
lernten schneller, endeten aber ebenfalls dort; Lektionen und tiefere Suche brachten nichts Messbares.
Frische Partien dagegen halfen. Der Schluss: Die Netze waren durch die \emph{Zahl der Partien} begrenzt,
nicht durch Größe oder Aufbau, solange das Netz nicht zu klein ist. Zum Vergleich: large hatte etwa
123\,000 Partien gesehen, AlphaZero im Schach 44 Millionen.

\paragraph{Zwei Tests für den finalen Lauf.} Mit je 5\,600\,s Rechenzeit wurden zwei Wege zu mehr Partien
pro Stunde verglichen: \emph{sims\_200} (200 statt 400 Simulationen pro Zug, etwa 1{,}7-mal so viele
Partien) und \emph{fast\_moves} (der Trick von KataGo~\cite{katago}: volle Suche nur auf einem Viertel der
Züge). Gegen \emph{more\_games} Runde 260 erreichte \emph{sims\_200} 56{,}8\,\% und bestand die vorher
gesetzte Schwelle von 53\,\%; \emph{fast\_moves} kam auf 51{,}8\,\%. Mehr Partien mit etwas flacherer
Suche schlugen also mehr Suche pro Partie.

\section{Der finale Lauf}

\paragraph{Einstellungen.} large wurde ab \emph{sims\_200} Runde 279 weitertrainiert: 1024 Partien pro
Runde, alle gleichzeitig, 200 Simulationen pro Zug, jede Stellung im Mittel zweimal gelernt, Speicher
1{,}5 Millionen Stellungen (etwa 28 Runden), Lernrate 0{,}0001 und für die letzten 15\,\% der Zeit
0{,}00003. Alle 8 Runden ein Check gegen \emph{large 240} (512 Partien, 100 Simulationen) und ein Snapshot.
Alle Partien wurden in kompakter Form archiviert, sodass sich das Training später wiederholen ließe.

\paragraph{Umfang.} Der Lauf endete bei Runde 1296: 1\,017 Runden, etwa 1\,041\,000 neue Partien, 37{,}3
Stunden Selbstspiel und Lernen (38{,}6 mit den Checks), also etwa achtmal so viele Partien wie der
ursprüngliche large-Lauf. Insgesamt hat das finale Netz gut 1{,}2 Millionen Partien gesehen.

\begin{figure}[ht]
\centering
\begin{tikzpicture}
\begin{axis}[
  width=0.96\textwidth, height=6.2cm,
  xlabel={Runde}, ylabel={Punkte gegen large 240 (\%)},
  xmin=270, xmax=1305, ymin=48, ymax=66,
  grid=major, grid style={gray!20},
  legend pos=south east, legend style={font=\small, draw=none, fill=white, fill opacity=0.8},
  tick label style={font=\small}, label style={font=\small},
]
\fill[blau!8] (axis cs:1120,48) rectangle (axis cs:1305,66);
\addplot[only marks, mark=*, mark size=1.1pt, color=grau] coordinates {
(280,51.0) (288,50.4) (296,51.1) (304,51.4) (312,52.6) (320,52.9) (328,54.6) (336,50.5) (344,53.8) (352,58.2) (360,53.0) (368,52.5) (376,54.4) (384,55.0) (392,54.9) (400,53.4) (408,56.1) (416,51.7) (424,54.6) (432,54.0) (440,56.3) (448,54.0) (456,56.2) (464,58.9) (472,52.0) (480,56.0) (488,54.2) (496,57.1) (504,58.1) (512,55.4) (520,55.1) (528,56.2) (536,55.6) (544,57.2) (552,56.7) (560,55.0) (568,56.5) (576,55.5) (584,57.5) (592,56.2) (600,58.0) (608,56.9) (616,59.8) (624,57.7) (632,55.3) (640,56.3) (648,58.0) (656,58.0) (664,56.6) (672,58.3) (680,54.6) (688,58.5) (696,57.0) (704,58.2) (712,57.7) (720,57.3) (728,56.9) (736,58.0) (744,58.1) (752,55.6) (760,60.3) (768,57.8) (776,59.6) (784,60.1) (792,60.1) (800,57.9) (808,57.0) (816,56.9) (824,57.9) (832,56.9) (840,58.3) (848,59.5) (856,58.1) (864,60.6) (872,58.6) (880,59.9) (888,60.6) (896,61.8) (904,63.4) (912,58.1) (920,60.6) (928,57.7) (936,60.3) (944,60.6) (952,60.4) (960,59.6) (968,58.9) (976,61.2) (984,60.4) (992,59.3) (1000,61.3) (1008,60.9) (1016,59.3) (1024,59.8) (1032,56.9) (1040,60.5) (1048,58.8) (1056,60.6) (1064,60.9) (1072,60.5) (1080,59.9) (1088,61.6) (1096,60.8) (1104,61.1) (1112,61.7) (1120,60.5) (1128,61.4) (1136,61.4) (1144,62.2) (1152,61.1) (1160,61.4) (1168,60.2) (1176,61.4) (1184,63.4) (1192,62.7) (1200,61.7) (1208,62.2) (1216,62.5) (1224,62.6) (1232,60.9) (1240,63.7) (1248,60.5) (1256,61.9) (1264,62.2) (1272,61.8) (1280,62.1) (1288,62.3) (1296,62.1)
};
\addlegendentry{Check (512 Partien)}
\addplot[thick, color=blau] coordinates {
(296,51.3) (348,53.7) (408,54.2) (468,55.6) (532,56.2) (596,57.3) (664,57.0) (728,57.4) (792,58.7) (856,58.7) (920,60.4) (984,60.3) (1048,59.7) (1100,61.0) (1148,61.2) (1208,62.5) (1268,61.9)
};
\addlegendentry{Mittel über Blöcke von 5--8 Checks}
\node[font=\small, anchor=north east, text=blau] at (axis cs:1300,65.5) {Lernrate 0{,}00003};
\end{axis}
\end{tikzpicture}
\caption{Checks des finalen Laufs gegen large 240 (100 Simulationen pro Zug). Ein einzelner Check schwankt
um etwa $\pm2$ Punkte. Der blaue Bereich ab Runde 1120 ist die kleine Lernrate.}
\label{fig:punkte}
\end{figure}

\begin{figure}[ht]
\centering
\begin{tikzpicture}
\begin{axis}[
  width=0.96\textwidth, height=5.4cm,
  xlabel={Runde}, ylabel={schlechtere Eröffnungen (von 256)},
  xmin=270, xmax=1305, ymin=10, ymax=80,
  grid=major, grid style={gray!20},
  tick label style={font=\small}, label style={font=\small},
]
\fill[blau!8] (axis cs:1120,10) rectangle (axis cs:1305,80);
\addplot[only marks, mark=*, mark size=1.1pt, color=grau] coordinates {
(280,67) (288,74) (296,60) (304,65) (312,59) (320,59) (328,56) (336,70) (344,48) (352,40) (360,67) (368,67) (376,54) (384,56) (392,42) (400,55) (408,47) (416,73) (424,53) (432,55) (440,42) (448,44) (456,51) (464,39) (472,60) (480,43) (488,59) (496,39) (504,45) (512,60) (520,51) (528,54) (536,49) (544,50) (552,47) (560,50) (568,53) (576,58) (584,44) (592,51) (600,51) (608,56) (616,38) (624,54) (632,58) (640,47) (648,46) (656,51) (664,53) (672,45) (680,57) (688,44) (696,53) (704,48) (712,35) (720,42) (728,47) (736,45) (744,49) (752,51) (760,41) (768,46) (776,46) (784,43) (792,44) (800,48) (808,49) (816,43) (824,41) (832,55) (840,41) (848,37) (856,41) (864,34) (872,43) (880,32) (888,41) (896,28) (904,20) (912,41) (920,30) (928,41) (936,35) (944,30) (952,24) (960,35) (968,42) (976,42) (984,28) (992,36) (1000,28) (1008,34) (1016,30) (1024,33) (1032,53) (1040,34) (1048,39) (1056,30) (1064,40) (1072,35) (1080,36) (1088,35) (1096,33) (1104,30) (1112,32) (1120,33) (1128,32) (1136,31) (1144,30) (1152,39) (1160,23) (1168,29) (1176,36) (1184,29) (1192,26) (1200,44) (1208,32) (1216,36) (1224,30) (1232,31) (1240,26) (1248,38) (1256,33) (1264,34) (1272,37) (1280,36) (1288,34) (1296,38)
};
\addplot[thick, color=rot] coordinates {
(296,65) (348,58) (408,54) (468,47) (532,51) (596,51) (664,50) (728,46) (792,45) (856,41) (920,33) (984,34) (1048,37) (1100,33) (1148,32) (1208,32) (1268,36)
};
\end{axis}
\end{tikzpicture}
\caption{Eröffnungen, auf denen das trainierte Netz in den Checks schlechter spielte als large 240 (von 256;
rote Linie: Mittel über Blöcke). Die Fehler gingen bis etwa Runde 950 deutlich zurück und blieben dann bei
etwa 32 bis 36.}
\label{fig:schlechter}
\end{figure}

\paragraph{Verlauf.} Abbildung~\ref{fig:punkte} zeigt die Punkte, Abbildung~\ref{fig:schlechter} die
schlechteren Eröffnungen. Nach einem Plateau um Runde 400 stiegen die Punkte stetig, etwa um einen Punkt
alle 60 Runden, später langsamer, bis auf rund 60\,\% um Runde 900. Die Zahl der schlechteren Eröffnungen
blieb zwischen Runde 440 und 650 bei etwa 50 stehen (der Gewinn kam damals aus gleichen Eröffnungen, die
zu besseren wurden) und fiel danach bis Runde 950 auf etwa 33. Die kleine Lernrate ab Runde 1120 hob die
Punkte noch einmal um etwa 2 auf rund 62\,\%, die schlechteren Eröffnungen aber nicht weiter. Auch die
Verluste und der Anteil Remis im Selbstspiel blieben ab etwa Runde 1000 flach.

\paragraph{Zwischenfälle.} Zwei Dinge kosteten Zeit. Erstens war Google Drive voll, als die Colab-Maschine
bei Runde 1262 unerwartet stoppte: Dateien werden zuerst lokal geschrieben und im Hintergrund hochgeladen,
und bei vollem Drive scheitert das ohne Meldung. Auf Drive lag nur das Netz von Runde 1167; die Runden
1168 bis 1260 (etwa 3{,}2 Stunden) gingen verloren und wurden neu gespielt. Zweitens schlief der Laptop
während der Endspiele mit zugeklapptem Deckel ein, was 3{,}4 Stunden kostete; verloren ging dabei nichts.

\section{Ergebnisse der Endspiele}

\subsection{Gegen die bisherigen Netze}

Alle Paare spielten die 300 verschiedenen Eröffnungen von beiden Seiten, also 600 Partien pro Paar und
Einstellung. Als Maßstab spielte zusätzlich large 240 gegen medium 240 auf denselben Eröffnungen.

\begin{table}[ht]
\centering
\footnotesize
\setlength{\tabcolsep}{5pt}
\begin{tabular}{llrrr}
\toprule
Paar & Sim. & Punkte & S -- U -- N & Eröffn.\ schlechter / gleich / besser \\
\midrule
final 1296 -- large 240 & 500  & 61{,}3\,\% & 170 -- 395 -- 35 & 22 / 144 / 134 \\
                        & 2000 & 58{,}2\,\% & 122 -- 455 -- 23 & 19 / 173 / 108 \\
                        & zus. & \textbf{59{,}8\,\%} & 292 -- 850 -- 58 & \textbf{41} / 317 / 242 \\
\midrule
final 1296 -- medium 240 & 500  & 64{,}0\,\% & 196 -- 376 -- 28 & 17 / 125 / 158 \\
                         & 2000 & 64{,}2\,\% & 188 -- 395 -- 17 & 9 / 133 / 158 \\
                         & zus. & \textbf{64{,}1\,\%} & 384 -- 771 -- 45 & \textbf{26} / 258 / 316 \\
\midrule
large 240 -- medium 240 (Maßstab) & 500  & 53{,}4\,\% & 117 -- 407 -- 76 & 57 / 154 / 89 \\
                                  & 2000 & 54{,}0\,\% & 98 -- 452 -- 50 & 35 / 179 / 86 \\
                                  & zus. & 53{,}7\,\% & 215 -- 859 -- 126 & 92 / 333 / 175 \\
\bottomrule
\end{tabular}
\caption{Direkte Partien. S -- U -- N: Siege, Unentschieden, Niederlagen des erstgenannten Netzes.
Eröffnungen nach der Eröffnungsregel aus dessen Sicht (je 300).}
\label{tab:netze}
\end{table}

Das finale Netz ist in beiden Bedenkzeiten klar das stärkste (Tabelle~\ref{tab:netze}). Mehr Simulationen
machen seinen Vorsprung kaum kleiner: Gegen large 240 sinkt die Punktzahl von 61{,}3 auf 58{,}2\,\%, weil
mehr Partien remis enden, aber die Zahl der schlechteren Eröffnungen bleibt mit 22 bzw.\ 19 von 300
niedrig. Der Maßstab zeigt, was das bedeutet: Der letzte große Schritt davor, large 240 gegen medium 240,
war noch auf 57 bzw.\ 35 Eröffnungen schlechter.

\subsection{Gegen Minimax}

Gegen Minimax wurden alle 102 verschiedenen Eröffnungen mit zwei Zügen von beiden Seiten gespielt. Die
ersten 53 davon sind die verschiedenen Eröffnungen, die das Netz \emph{medium 132} im ersten großen
Minimax-Test gespielt hatte; auf ihnen lassen sich beide direkt vergleichen.

\begin{table}[ht]
\centering
\footnotesize
\setlength{\tabcolsep}{4pt}
\begin{tabular}{llrrrr}
\toprule
Sim. & Tiefe & final 1296 & schlechter / gleich / besser & final, 53 Eröffn. & medium 132, 53 Eröffn. \\
\midrule
500  & 13 & 88{,}0\,\% (163--33--8)  & 3 / 6 / 93   & 90{,}1\,\%, 0 / 2 / 51 & 73{,}1\,\%, 3 / 17 / 33 \\
500  & 15 & 79{,}9\,\% (133--60--11) & 5 / 16 / 81  & 77{,}8\,\%, 3 / 9 / 41 & 62{,}3\,\%, 10 / 13 / 30 \\
2000 & 13 & 94{,}1\,\% (182--20--2)  & 0 / 2 / 100  & 95{,}8\,\%, 0 / 1 / 52 & 84{,}9\,\%, 0 / 6 / 47 \\
2000 & 15 & 88{,}2\,\% (159--42--3)  & 1 / 7 / 94   & 89{,}2\,\%, 0 / 4 / 49 & 75{,}5\,\%, 3 / 13 / 37 \\
\midrule
\multicolumn{2}{l}{zusammen} & \textbf{87{,}6\,\%} (637--155--24) & \textbf{9} / 31 / 368 & & \\
\bottomrule
\end{tabular}
\caption{Partien gegen Minimax, 204 pro Zeile (auf den 53 gemeinsamen Eröffnungen 106).}
\label{tab:minimax}
\end{table}

Gegen Minimax verliert das finale Netz mit 2000 Simulationen nur noch 2 bzw.\ 3 von 204 Partien und ist
selbst gegen Tiefe 15 nur auf einer von 102 Eröffnungen schlechter (Tabelle~\ref{tab:minimax}). Auf
denselben Eröffnungen holt es 11 bis 17 Punkte mehr als medium 132 und verliert etwa dreimal seltener.

\section{Diskussion}

\paragraph{Was geholfen hat.} Der wichtigste Befund des Projekts ist schlicht: \emph{Partien}. Größe,
Aufbau, zusätzliche Lektionen und tiefere Suche brachten bei gleicher Datenmenge kaum etwas; frische
Partien halfen in jedem Versuch. Konsequent war deshalb, die verbleibende Rechenzeit fast vollständig in
einen langen Lauf mit möglichst vielen Partien pro Stunde zu stecken und dafür sogar die Suche pro Zug zu
halbieren. Die Endspiele bestätigen diese Wahl: Das Netz wurde nicht nur in den Punkten, sondern auch nach
der strengeren Eröffnungsregel deutlich besser.

\paragraph{Wo es aufhört.} Die Checks flachten ab etwa Runde 1000 ab, und die schlechteren Eröffnungen
fielen auch mit kleinerer Lernrate nicht weiter unter etwa 32. Verluste und Remisanteil im Selbstspiel
waren ebenfalls flach. Das deutet darauf hin, dass das Netz bei dieser Größe und diesen Einstellungen nahe
an seinem Ende ist. Weitere Runden hätten vermutlich wenig gebracht; ein nächster Schritt bräuchte eher
ein größeres Netz, das von den vielen Partien profitieren kann, oder bessere Partien pro Rechenzeit.

\paragraph{Grenzen der Messung.} Einige Einschränkungen gehören zu den Ergebnissen:
\begin{itemize}
\item \emph{Zufall.} 600 Partien sind auf etwa $\pm1{,}5$ Punkte genau, 204 auf etwa $\pm2{,}5$ (geschätzt).
Die großen Unterschiede in den Tabellen liegen deutlich darüber, kleine Unterschiede zwischen Zeilen nicht.
\item \emph{Der Minimax-Vergleich.} medium 132 ist kleiner und viel kürzer trainiert. Der Vergleich zeigt den
Fortschritt des ganzen Projekts, nicht den des finalen Laufs allein. Ein direkter Maßstab wäre large 240
gegen Minimax auf denselben Eröffnungen; das wurde nicht gespielt.
\item \emph{Checks bei 100 Simulationen} überschätzen Unterschiede gegenüber längerer Bedenkzeit. In den
Endspielen blieb der Vorsprung bei 500 Simulationen aber gleich groß wie in den Checks (61 bis 62\,\%) und
sank bei 2000 nur wenig.
\item \emph{Kein starker externer Gegner.} Der einzige Vergleich mit einem fremden Programm (das Netz von
uttt.ai, mit unserer Suche) war mit medium ausgeglichen. Wie stark das finale Netz im Vergleich zu den
besten Programmen ist, bleibt offen.
\end{itemize}

\paragraph{Symmetrie.} Das Netz ist nicht genau symmetrisch: Gibt man ihm dieselbe Stellung in allen acht
Drehungen und Spiegelungen, unterscheiden sich die Zugverteilungen im Mittel um 6{,}7\,\% und die
Bewertungen um 0{,}09 (auf der Skala von $-1$ bis $+1$). Das finale Netz ist dabei symmetrischer als large
240 (7{,}8\,\%) und medium 240 (9{,}1\,\%). Die Symmetrie wird nur gelernt, weil jede Stichprobe zufällig
gedreht wird, nicht durch den Aufbau erzwungen. Ein billiger, nicht getesteter Gewinn wäre, die Stellung
bei jeder Netzabfrage in der Suche zufällig zu drehen, wie AlphaGo Zero~\cite{alphagozero} und
KataGo~\cite{katago} es tun: Die kleinen Fehler mitteln sich dann über die vielen Simulationen heraus.

\paragraph{Was offen bleibt.}
\begin{itemize}
\item Gewichtsmittelung der letzten Snapshots (sie liegen nur auf Drive).
\item Zufällige Drehung in der Suche (siehe oben), testbar ohne neues Training.
\item large 240 gegen Minimax als direkter Maßstab.
\item Ein größeres Netz, trainiert aus dem Archiv der rund einer Million gespeicherten Partien. Für large
kostete ein solcher Lauf auf der L4 etwa drei Stunden, für ein größeres Netz entsprechend mehr; er würde zeigen, ob Größe bei dieser Datenmenge nun hilft.
\end{itemize}

\paragraph{Zur Arbeitsweise.} Hilfreich war, Entscheidungsregeln \emph{vor} einem Lauf festzulegen
(Schwellen von 53\,\% bzw.\ 48\,\%, Stopp-Regeln) und Netze nach direkten Partien statt nach Verlusten oder
Checks gegen ein drittes Netz zu beurteilen. Das Linien-Netz etwa führte in den Checks gegen medium und
verlor trotzdem direkt gegen large. Ebenso nützlich war die Eröffnungsregel: Sie zeigte, dass der
Fortschritt im finalen Lauf zeitweise nur aus gleichen Eröffnungen kam, die zu besseren wurden, und erst
später aus weniger Fehlern.

\section{Fazit}

Das Projekt hat sein Ziel erreicht: Das finale Netz ist ein reines, selbst gelerntes Netz, das deutlich
stärker spielt als alle Vorgänger und vor allem \emph{seltener schlechter} als sie. Gegen large 240 ist es
auf nur 41 von 600 Eröffnungen schlechter (der vorherige Schritt: 92), gegen Minimax der Tiefe 15 mit
2000 Simulationen auf einer von 102. Der Weg dorthin war nicht ein besserer Aufbau, sondern die Einsicht
aus vielen gleich endenden Versuchen, dass die Netze vor allem mehr Partien brauchten, und die
Entscheidung, die Rechenzeit darauf zu konzentrieren. Zugleich zeigen die flachen Kurven am Ende, dass
dieses Netz mit diesen Einstellungen weitgehend ausgeschöpft ist. Wer weitergehen will, findet im Archiv
der gespeicherten Partien die Grundlage für einen nächsten, größeren Schritt.

\begin{thebibliography}{9}
\bibitem{alphazero} D.~Silver et~al.: \emph{A general reinforcement learning algorithm that masters chess,
shogi, and Go through self-play.} Science 362 (2018).
\bibitem{alphagozero} D.~Silver et~al.: \emph{Mastering the game of Go without human knowledge.} Nature 550
(2017).
\bibitem{katago} D.~J.~Wu: \emph{Accelerating Self-Play Learning in Go.} arXiv:1902.10565 (2019).
\end{thebibliography}

\end{document}
